

\section{Clifford-free Depth Overhead.}
\label{section:cliffordfree}
Our next objective is to modify the construction to introduce no depth overhead for the Clifford gates. As already seen, the $\phi$-transpose followed by transversal Hadamard exchanges $X_{S^{\prime}} \leftrightarrow Z_{S/S^{\prime}}$. Therefore, when $d' = d - d'$ and with the help of permutation over the generators, a logical Hadamard has a fold-transversal realization. Unfortunately, the natural realization does not respect the size function. Yet, for a slightly different function, which we think of as the diameter function, we will be able to set reasonable bounds on its growth. We start by defining the diameter of a finite chain.

\newcommand{\Sizestar}[1]{ \mathbf{Size}^{\star}( #1 )  }
Next, we introduce the realization of a logical Hadamard.
For a fixed subset of generators $ J \subset S $, which is used to define the encoded logical bit, we define the permutation that permutes $ J $ into $ S/J $ and vice versa. Extending this permutation, first to an automorphism over the Toric elements, and then to a permutation over the Toric faces, gives a map that keeps the stabilizers but sends $ X_{J} \leftrightarrow X_{S/J} $.
\begin{definition}
Let $\psi \colon S \rightarrow S$ be a permutation on the generators in $S$, so $\psi$ has a natural extension to an automorphism on $\mathcal{G}$\footnote{By replacing each of the generators in a product by its image according to $\psi$.}. Extend $\psi$ to act over the faces of the $(d,d^{\prime})$-Toric as follows:
  \begin{equation*}
    \begin{split}
      \psi( \rotfvs{}) \colonequals \rotfv{\psi(v)}{\psi(S^{\prime})}
    \end{split}
  \end{equation*}
  We name the map a $\psi$-permutation. 
\end{definition}

\begin{claim}
  The $\psi$-permutation sends $X$-checks to $X$-checks and $Z$-checks to $Z$-checks. Moreover, assume that $d^{\prime}=d/2$. Then, for a generator subset $J\subset S$ such that $d^{\prime} = |J| = |S/J|$ and $S/J = \psi(J)$, we have that the application of $\psi$ realizes the exchange $x_{J} \leftrightarrow x_{S/J}$ and $z_{J} \leftrightarrow z_{S/J}$. 
\end{claim}

\begin{proof}
Consider the action of $\psi$ on a $d'+1$-dimensional face $\rotfvs{}$:
  \begin{equation*}
    \begin{split}
      \psi \left( \partial^{-} \rotfvs{} \right) &= \psi \left( \sum_{g\in S^{\prime}} \rotfv{v}{S^{\prime}/g} + \rotfv{gv}{S^{\prime}/g}   \right) \\
      &= \sum_{g\in S^{\prime}} \rotfv{ \psi(v)}{\psi(S^{\prime})/\psi(g)} + \rotfv{\psi(g)\psi(v)}{\psi(S^{\prime})/\psi(g)}  \\
      &= \sum_{g\in \psi(S^{\prime}) } \rotfv{ \psi(v)}{\psi(S^{\prime})/g} + \rotfv{g\psi(v)}{\psi(S^{\prime})/g}  \\
      &= \partial^{-}\rotfv{\psi(v)}{\psi(S^{\prime})}
    \end{split}
  \end{equation*}
Since $\psi$ is a permutation over $S$, it follows that $|\psi(S^{\prime})| = |S^{\prime}| = d^{\prime} + 1$; namely, $\partial^{-}\rotfv{\psi(v)}{\psi(S^{\prime})}$ is a $Z$-check. The same calculation shows that for $S^{\prime}$ at size $|S^{\prime}|=d^{\prime}-1$, we have $\psi \left( \partial^{+} \rotfvs{} \right) = \partial^{+}\rotfv{\psi(v)}{\psi(S^{\prime})}$. 

Now, consider $J\subset S$ such $|J|=d^{\prime}$ and $\psi(J) = S/J$. Then: 
\begin{equation*}
  \begin{split}
    \psi(x_{J}) =  \psi \left( \sum_{v \in \expval*{J}} \rotfv{v}{J} \right) = \sum_{v \in \expval*{J}} \rotfv{\psi(v)}{\psi(J)} = \sum_{v \in \expval*{S/J}} \rotfv{v}{S/J} 
  \end{split}
\end{equation*}
Where we used the fact that $\psi$ sends the group generated by $J$ to the group generated by the complement generator set. Since $\psi$ is a permutation, $\psi(S/J) = J$, and therefore $\psi(x_{S/J}) = x_J$. To get the exchange $z_{J} \leftrightarrow z_{S/J}$, we repeat the above, where the root vertices in the sum are the vertices in the subgroup generated by the complement of the subset of generators spanning the face.
\end{proof}
The corollary is that in the \textit{nice encoding} (\Cref{def:niceenc}), the logical Hadamard has a fault-transversal implementation.
\begin{corollary}
  \label{corollary:hadamard}
  For $ d' = d - d' = d/2 $, one can fix a generator subset $ S' \subset S $ of size $|S'| = d^{\prime}$, and realize a Hadamard gate that exchanges $ X_{S'} \leftrightarrow Z_{S'} $.
\end{corollary}
\begin{proof}
  The realization is obtained by the concatenation $H^{\otimes n} \circ \phi \circ \psi$.
\end{proof}
For a given finite chain (that might support an error), the following two claims set a bound on the increase in its diameter by the application of logical Hadamard. Specifically, the diameter grows by at most $2 + d^{\prime}$, and therefore, in the fault-tolerant construction, we will later set $d^{\prime} + 2$ layers of sweep-cycle after any Hadamard application.
\begin{claim}
  \label{claim:invpar}
  Let $z$ be a connected $d^{\prime}$-dimensional chain. Assume that both $z$ and $\psi(z)$ are finite. Then we have:
  \begin{equation*}
    \begin{split}
      \max_{v \in z} L_{d+1}(v)  = \max_{v \in \psi(z)} L_{d+1}(v)  \, \, \, \text{   and   }  \, \, \, \max_{g\in S}\max_{v \in z} L_{d+1}(v)  = \max_{g\in S}\max_{v \in \psi(z)} L_{g}(v)
    \end{split}
  \end{equation*}
\end{claim}
\begin{proof}
Follows directly from the fact that summation and the max-function are invariant to permutations on the set they act on.
\end{proof}

\begin{remark}
\inb{\textbf{Change Here.}}  Notice that although the statement in \Cref{claim:invpar} is correct for the size function ($\Size{\cdot}$), it fails in the language of slices. For example, consider a slice that spreads over two registers, a subset of generators $S^{\prime} \subset S$, and a generator $g \in S^{\prime}$, such that both registers have the error of $X_{z}$, where:
  \begin{equation*}
    \begin{split}
      z \colonequals \rotfv{v}{S^{\prime}}+ \rotfv{gv}{S^{\prime}} + \rotfv{g^{2}v}{S^{\prime}} + ,.., + \rotfv{g^{10}v}{S^{\prime}}
    \end{split}
  \end{equation*}
Now, consider permuting by $\psi$-permutation only the qubits in the first register. Clearly, the diameter of the slice increases.
\end{remark}

\begin{claim}
  \label{claim:diamexp}
  Consider a $d^{\prime}$-dimensional chain $z$ such both $z$ and $\phi(z)$ are finite, then it holds that: 
      \begin{equation*}
        \begin{split}
          \bigl| \Size{ z  } - \Size{ \phi(z)  } \bigr| < 2d + d^{\prime}
        \end{split}
      \end{equation*}
\end{claim}
\begin{proof}
  We start by showing that: 
\begin{equation*}
  \begin{split}
    \max_{v \in \phi(z)}L_{d+1}(v) = \max_{v \in z }L_{d+1}(v) - d^{\prime}
  \end{split}
\end{equation*}
Let us consider the action of $\phi$ on a rooted chain. Let $y$ be a rooted chain at vertex $v$, namely there are faces $A = \{ \rotfvs{} \}$, such that:
\begin{equation*}
  \begin{split}
    y = \sum_{ \rotfvs{} \in A}\rotfvs{}
  \end{split}
\end{equation*}
Thus, 
\begin{equation*}
  \begin{split}
    \phi(y) = \sum_{ \rotfvs{} \in A}\rotfv{S^{\prime}v}{S/S^{\prime}}
  \end{split}
\end{equation*}
So if $v$ achieves the maximal $L_{d+1}$-value in $z$, then the vertices that get the maximal value in $\phi(z)$ are of the form $S^{\prime}v$, for some $d^{\prime}$-size subset of generators $S^{\prime} \subset S$. Therefore, the following equality holds:
\begin{equation*}
  \begin{split}
    \max_{v \in \phi(z)}L_{d+1}(v) = \max_{v \in z }L_{d+1}(v) - d^{\prime}
  \end{split}
\end{equation*}
So it's left to show that for any generator $g \in S$ it holds that: 
\begin{equation*}
  \begin{split}
     \bigl| \max_{v \in \phi(z)}L_{g}(v) -  \max_{v \in z }L_{g}(v) \bigr| \le 1
  \end{split}
\end{equation*}
Let $v$ be a vertex that is a root of a face in $\phi(z)$. We argue that $v$ is also in $z$. Consider a face $\rotfvs{}$ in $\phi(z)$, then $\phi^{-1}(\rotfv{v}{S^{\prime}}) = \rotfv{(S/S^{\prime})^{-1} v}{S/S^{\prime}}$ is in $z$, and therefore $(S/S^{\prime})(S/S^{\prime})^{-1} v = v$ is in $z$. By the same arguments, we also have that if $v$ is a root of a face in $z$, then the vertex $S v$ is in $\phi(z)$.


Now, fix a generator $g \in S$ and let $u$ and $w$ be the vertices in $z$ and $\phi(z)$ that achieve the maximum for $L_{g}$. Denote by $u_{r}$ and $w_{r}$ the vertices which are the roots of the faces containing $u$ and $w$, and observe that $L_{g}(w) \le L_{g}(w_{r})+1$ for any $g \in S$. We have that:
    \begin{equation*}
      \begin{split}
        \max_{v \in z} L_{i}(v) &  \le   L_{g}(S u_{r}) \le  \max_{v \in \phi(z)} L_{g}(v)  \\
        \max_{v \in \phi(z)} L_{g}(v) & \le  L_{g}(w_{r})  + 1 \le  \max_{v \in z} L_{g}(v) + 1 
      \end{split}
    \end{equation*}
So, the inequality holds.
\end{proof}

\begin{definition}
  For a $(2d^{\prime},d^{\prime})$-Toric code, we define the realization of Hadamard gate by $ H^{\otimes (\cdot) } \circ \phi \circ \psi$ followed by $d^{\prime} + 2d$ times sweep-cycle (See also \Cref{corollary:hadamard}). 
\end{definition}

Having a realization of logical Hadamard, we would like to extend the Toric encoded circuits to use it instead of gate teleportation. In addition, we define the Clifford-free Toom's contour to be the space-time graph that describes a noisy running of the extended encoding circuit. The definition of the Clifford-free Toric contour presents two main differences. First, the vertices are not partitioned into $X$ and $Z$ types; instead, a vertex would belong to $V_{\mathcal{H}}$ if it supports a face in the reduced Pauli-path regardless of its type. Second, the arrows connect vertices that are at a 'time-distance' of $d^{\prime}+2d +1$ since that is the time the computation layer lasts.

\begin{definition}[Clifford-free Toric Encoded Circuit]
We extend the Toric Encoded definition to the Clifford-free Toric encoded circuit by requiring that $d^{\prime}=d-d^{\prime}=d/2$, replacing the realization of the Hadamard gate with the fold-transversal realization instead of Gate-teleportation. In addition, we enforce that the depth of the logical layers be exactly $d^{\prime}+2d+1$ by padding with identity\footnote{And, of course, in any practical implementation, it will be better to pad with Sweep-cycles. We do it for the convenience of the adaptation.}.
\end{definition}

\begin{definition}[Clifford-free Toom's Contour]
    For a Clifford-free Toric encoded quantum circuit $C$, set of faults $\mathcal{E}$ and a reduced Pauli-path $P_{1},..$, let us denote the base graph of $C[\mathcal{E}]$ by $G_{o}$. We define the \textit{Clifford-free Toom's contour graph} $\mathcal{H} = ( V_{\mathcal{H}}, E_{\mathcal{H}})$ as follows: 
    \begin{enumerate}
      \item The vertices $V_{\mathcal{H}}$ correspond to vertices of faces supporting the reduced Pauli-path. 
        \begin{equation*}
          \begin{split}
            V_{\mathcal{H}} \colonequals & \big\{ v_{t} : v_{t}=(v,t) \in V_{o} \\
            & \, \, \, \text{ and there exists a face } f \in \text{ supp } P_{t} \text{ such } v \in f \big\} \\
          \end{split}
        \end{equation*}
      \item The edges $E_{\mathcal{H}}$ is partitioned into two types $A_{\mathcal{H}}$ and $F_{\mathcal{H}}$, where for any logical gate $g$ which applied in time $t$ in $C[\mathcal{E}]^{\infty} $ we add edge to $A_{\mathcal{H}}$ as follow: 
        \begin{enumerate}
    \item Any logical gate that is not Hadamard connects vertices at the initial time of the layer with vertices at the final time of the layer if they share a face in the Pauli-path. Since the padding is done by identity, it's equivalent to connect $(v,t)$ with $(u,t+d^{\prime}+2d+1)$ by an arrow if both $(v,t)$ and $(u,t+1)$ are in $V_{\mathcal{H}}$, and there are faces $f_{1}, f_{2}$ (which might be the same) and a physical gate $g^{\prime}$ acting at time $t$ such that $v \in f_{1} \in \text{Loc}(g^{\prime})$ and $u \in f_{2} \in \text{Loc}(g^{\prime})$.
    \item If the logical gate is a Hadamard, then we connect $(v,t)$ with $(u,t+d^{\prime}+2d+1)$ if there are vertices $(v_{0},t),(v_{1},t+1),(v_{2},t+2),\ldots,(v_{d^{\prime}+2d},t+d^{\prime}+2d),(v_{d^{\prime}+2d+1},t+d^{\prime}+2d+1) \in V_{\mathcal{H}}$ such that $v_{0}=v$, $v_{d^{\prime}+2d+1}=u$, and for any $i \in [0, d^{\prime}+2d+1]$, there is a physical gate $g^{\prime}$ acting at time $t + i$ and a face $f \in \text{Loc}(g^{\prime})$ containing both $v_{i}$ and $v_{i+1}$.
    \item If $g$ is a fault Pauli, we add:
            \begin{equation*}
              \begin{split}
                & \bigl\{  \{  \tilde{u} = (u,t), \tilde{v} = (v,t+d^{\prime}+2d+1)\} : \text{ if } \tilde{v},\tilde{u} \in V_{\mathcal{H}}  \text{, there exists a face } f \text{ such }  u,v \in f \text{ , }  \text{ and }  \\ 
                &  \, \, \, \text { for any face  } f^{\prime} \text{ such }  u,v \in f^{\prime} \text{ it holds that  } f^{\prime} \notin  \text{ Supp } g  \bigr\} 
              \end{split}
            \end{equation*}
        \end{enumerate}
        To construct $F_{\mathcal{H}}$, we connect any pair of vertices $(v,t)$, $(u,t)$ if there exists a vertex $(w,t+d^{\prime}+2d + 1)$ such both of them are connected to it trough $A_{\mathcal{H}}$. Namely: 
        \begin{equation*}
          \begin{split}
            F_{\mathcal{H}} &= \big\{ e = \{u,v\} : \text{ exists } w \in V_{\mathcal{H}} \text{ such } \{v,w\},\{u,w\} \in A_{\mathcal{H}} \big\} 
          \end{split}
        \end{equation*}
    \end{enumerate}
    We call $A_{\mathcal{H}}$ and $F_{\mathcal{H}}$ the \textit{arrows} and \textit{forks} edges respectively. 
\end{definition}


%\begin{definition}
%  Denote the quantum registers by $R_{1},R_{2},.., R_{r}$. Let $K$ be a slice, let $K_{1},..,K_{i},..,K_{r}$ be the restrictions of $K$ to the registers. And let us decompose each restriction to register $K_{i}$ into connected component $K_{i}^{1},..,K_{i}^{m_{i}}$, So: 
%  \begin{equation*}
%    \begin{split}
%      K =  \bigcup_{i=1}^{r}\bigcup K_{i}^{j}
%    \end{split}
%  \end{equation*}
% For a $r$-length binary string $x \in \{0,1\}^{r}$, a $\psi_{x}$-configuration of $K$ is 
% \begin{equation*}
%   \begin{split}
%     \psi_{x}(K) \colonequals \bigcup_{x_{i}=0} K_{i} \bigcup_{x_{i}=1} \psi(K_{i})
%   \end{split}
% \end{equation*}
% Where 
% \begin{equation*}
%   \begin{split}
%     \psi(K_{R_{i}}) = \big\{ \psi(v) : v \in R_{i} \big\}
%   \end{split}
% \end{equation*}
%\end{definition}

\begin{definition}
Let $A$ be either a slice or a subslice. We define the size-star function to assign $A$:
\begin{equation*}
  \begin{split}
    \Sizestar{A} \colonequals \Size{ A \cup \psi(A)  } 
  \end{split}
\end{equation*}
\end{definition}


\inb{ \textbf{ For both the repelacment and the intermidiate arguments the trick is to show that $  \Size{ \cdot } \le \Sizestar{  \cdot  } \le d \Size{ } $.   } }

\begin{claim}
  \begin{equation*}
    \begin{split}
     \Size{ \cdot } \le \Sizestar{  \cdot  } \le d \Size{ } 
    \end{split}
  \end{equation*}
\end{claim}
\begin{proof}
  The first ineqaulity follows by the contianment $A \subset A \cup \psi(A)$.  For the second: 
\end{proof}




\begin{claim}
  The replacement argument holds when substituting $\Size{ }$ with $\Sizestar{ }$. 
\end{claim}

\begin{claim}
   Invariant of $\Sizestar{}$ under $\psi$. 
\end{claim}
\begin{proof}
  It's clear.  
\end{proof}

\begin{claim}
  \label{claim:modifitringinq}
Let $A, B$ be slices, such that $A \cap B \neq \emptyset$. Then:
  \begin{equation*}
    \begin{split}
      \Sizestar{A\cup B} \le \Sizestar{A} + \Sizestar{B} 
    \end{split}
  \end{equation*}
\end{claim}
\begin{proof}
First, notice that $\psi(A \cup B) = \psi(A) \cup \psi(B)$. Thus:
  \begin{equation*}
    \begin{split}
    \Sizestar{A\cup B} & =  \Size{ (A \cup \psi(A)) \cup  (B \cup \psi(B)) } \le   \Size{ (A \cup \psi(A)) } +  \Size{(B \cup \psi(B)) } \\
      & = \Sizestar{A} + \Sizestar{B} 
    \end{split}
  \end{equation*}
Where, in the first transition, we used \Cref{claim:tringinq}.
\end{proof}


\inb{ \textbf{ We need after CX perform two sweep-cycels. In the proof, decompise by $K_{i}$ and not by $K_{i}^{j}$ as different connected component might merge. } }




\begin{claim}
  The intermediate argument 'holds' when substituting $\Size{ }$ with $\Sizestar{ }$. 
\end{claim}
\begin{proof}
  Imagine that any $\psi$-permutation is grouped together into a single step. Assume that at step $\tau$ there is a subslice at size-star greater than or equal to $n$. If at $\tau$ we permute by $\psi$, then by the invariant of $\Sizestar{}$ under $\psi$, we also have that $\tau -1$ has a subslice at size-star greater than or equal to $n$. Thus, we can assume that the gate applied at turn $\tau$ is not a part of $\psi$-permutation. So, $u$ cause a merging. 
  Now, let $J_{1},..,J_{2k}$ be a connected set. If $J_{1}$ is connected to $J_{2}$ then they share an $L_{d+1}$ value (their intersection point achieve it). Therefore $\psi_{y}(J_{1})$ and $\psi_{y}(J_{2})$ share an $L_{d+1}$ value. Therefore: 
  \begin{equation*}
    \begin{split}
      \Sizestar{ J_{1} \cup J_{2} ,..    } \le \Size{ J_{1} \cup J_{2} ,..    } \le \sum \Size{ J_{i} } 
    \end{split}
  \end{equation*}
  Thus:
\begin{equation*}
  \begin{split}
    & \Sizestar{K}  \le \Delta d^{2} ( \Size{ \psi_{y}(K^{\prime}) } +  \Size{\text{fork}} ) \\
    \Rightarrow &  \Sizestar{ K^{\prime} } \ge  \psi_{y}(\Size{K^{\prime}}) \ge \frac{1}{\Delta d^{2}} \Sizestar{K} - \Size{ \text{fork} }
   \end{split}
\end{equation*}
\end{proof}


\inb{ \textbf{ The trick is to treat forks differently (like either apply $\psi$ over the two ends of a fork, or not apply at all). And we can do that because we still covers all the vertices in $\psi(K)$. } }

\begin{claim}
  The minimal explantation argument 'holds' when substituting $\Size{ }$ with $\Sizestar{ }$. 
\end{claim}

\begin{proof}
  In the induction step, 
  \begin{equation*}
    \begin{split}
      A & \le d+1 + \sum_{i=0}^{k}\alpha(m_{i} - \mathcal{X}_{r-1}) - \beta s^{\star}_{i} \\
          &\le  d+1 + \sum_{i=0}^{k}\alpha(m_{i} - \mathcal{X}_{r-1}) - \beta s_{i} \\
          & \le  d+1 + \left(\sum_{i=0}^{k}\alpha(m_{i} - \mathcal{X}_{r-1}) \right) - s^{\star}   
    \end{split}
  \end{equation*}

\end{proof}




With adjustments being made, we present a threshold theorem with a Clifford-free depth overhead.
  \begin{theorem}
    \label{theorem:cliffordfree}
    There exists a constant $ p_{0} \in (0,1) $ such that for any local stochastic noise channel $\mathcal{E}$ with rate $ p < p_{0} $, and a Clifford+$\mathbf{T}$ quantum circuit $ C $ with depth $\mu$, $\mathbf{T}$-depth $\mu_{\mathbf{T}}$ and width $ w $, there exists a Clifford-free Toric encoded circuit $\tilde{C}$, with depth $ \Theta(\mu) + \Theta( \mu_{\mathbf{T}} n^{5}) $ and width $\Theta((w + \mu) n^{8} \text{Polylog}(n) )$ that simulates $ C $ with probability:
\begin{equation*}
  \begin{split}
    1 -   O \left( \text{poly}(w,\mu,n) \right)  p^{\Omega(n)} 
  \end{split}
\end{equation*}
  \end{theorem} 
  \paragraph{Sketch of adjusting the proof of \Cref{theorem:main} to Clifford-free.}    
First, we change the definition of infinite slice/subslice/error to support the diameter function. Namely, a slice $K$ would be called infinite if $\Size{K} \ge n$. Then, the proofs of \Cref{claim:intersyn}, \Cref{claim:decodestopnoise}, and \Cref{claim:replace} can be easily modified to support an 'intermediate claim' and a 'replacement claim'. The first states that having a step that admits an infinite slice implies that a previous step admits a slice (according to \Cref{claim:rereduced}) with a diameter that is at least a constant fraction of the side of the Toric. The second states that the probability that an ideal-cleaning procedure fails is lower than the occurrence of an infinite slice in the Toom's contour obtained by replacing the ideal procedure with a noisy one.
  .
Next, we reformulate the explanation definition such that in the bounds on arrows, the diameter substitutes the size. Using \Cref{claim:prevslicesdiam} combined with the fact that the Hadamard realization ends with $d^{\prime}+2d$ sweep-cycles, one can modify \Cref{claim:edgesin} to show that assuming that up to a step $\tau$ no infinite slice was introduced, then any slice in the Clifford-free Toom's contour has an explanation. In particular, it's possible to bypass the expansion in diameter that might be introduced by transposing registers by $\phi$-transpose (\Cref{claim:diamexp}). From there, up to constants\footnote{For example, the maximal length of a fork is now $2 \cdot (d^{\prime}+2d+1) \cdot d$.}, the proof is identical to the proof of \Cref{lemma:keylemma}.


